\documentclass[12pt]{article}
\usepackage[utf8]{inputenc}
\usepackage[T1]{fontenc}
\usepackage{hyperref}
\title{Aloápam Zapotec Language Resources for K-12 Curriculum Development}
\author{Compiled by Claude Sonnet 5.5 (claude-sonnet-5-5) for Andr\'as Kornai}
\date{}
\begin{document}
\maketitle

\section{Introduction}
Aloap\'am Zapotec (ISO 639-3 zaq; Glottocode aloa1235) is an Otomanguean language of the Sierra Norte of Oaxaca, Mexico, with about 3,400 speakers, centred on San Miguel Aloap\'am and San Isidro Aloap\'am in the Ixtl\'an District. This report is based on a small number of web searches. Very little material was found, and nothing suitable as classroom curriculum material was located.

\section{Primary Educational Resources}
No primers, textbooks or school programmes for Aloap\'am Zapotec were found in the search results.

\section{Linguistic and Reference Materials}
\begin{itemize}
\item Omniglot, Aloap\'am Zapotec: a short overview, alphabet chart and sample text; it notes a Latin-alphabet orthography. \url{https://omniglot.com/writing/zapotec-aloapam.htm}
\item Glottolog entry for Aloap\'am Zapotec (language classification and bibliography pointers). \url{https://glottolog.org/resource/languoid/id/aloa1235}
\end{itemize}
No curated dictionary or grammar was found. Glosbe pages exist but are assembled from parallel corpora and are unverified, so they are not listed as a reference.

\section{Religious and Community Resources}
Omniglot reports that parts of the Bible have been translated into the language. Global Recordings Network lists Scripture resources (Scripture Earth, Faith Comes By Hearing) on its page for the language: \url{https://globalrecordings.net/en/language/zaq}

\section{Digital and Audio Resources}
The Global Recordings Network page for zaq lists audio material, though it notes that some recordings may have been withdrawn or are being remade. \url{https://globalrecordings.net/en/language/zaq}

\section{Curriculum Development Context}
The evidence found is thin: no school materials, no dictionary, and no community organisation were identified. Any curriculum would need to start from fieldwork with speakers and from the Latin-alphabet orthography used in the Scripture translation. Further searching in Spanish and in Mexican state education sources (for example the Oaxaca indigenous education institute) is recommended.

\end{document}